\documentclass[12pt]{article}
\usepackage{amsmath}
\usepackage{amssymb}  % Provides \mathbb and other symbols
\usepackage{geometry}
\usepackage{graphicx}
\usepackage{subcaption}
\usepackage{enumitem}
\usepackage{listings}
\usepackage{hyperref}
\usepackage{ulem}
\usepackage{xcolor}
\usepackage{amsthm}
\usepackage{tcolorbox}  % For colored boxes
\usepackage[table]{xcolor} % For coloring table cells
\usepackage{array}         % For better table alignment
\usepackage{mdframed}
\usepackage{cancel}
\usepackage{amsmath,amssymb}

% Set the top margin

\geometry{top=1in} % You can adjust the measurement as needed


\begin{document}
\begin{titlepage}
\centering
\vspace*{\fill} % Adds vertical space to center the title block vertically
{\Huge Topology} \\[0.5cm] % Title 2
{\Huge Home Work 4}\\[1.5cm] % Title 3
{\Large Student: 26-712-224}\\[0.5cm] % Student ID
{\Large \today} % Date, \today adds today's date
\vspace*{\fill}
\end{titlepage}

\newpage
\noindent
\noindent\textbf{Exercise 1.} Let $(X,\mathcal{U})$ be a topological space.
\begin{enumerate}
\item[(a)] Let $C\subseteq X$ be connected. Prove that its closure $\overline{C}$ is connected.
\item[(b)] Let $C_1,C_2\subseteq X$ be connected and assume that $C_1\cap C_2\neq\emptyset$. Prove that $C_1\cup C_2$ is connected.
\item[(c)] Let $p\in X$, and consider the connected component
\begin{align*}
C(p):=\bigcup_{\substack{C\subseteq X\ \text{connected}\\p\in C}} C \tag{1.1}
\end{align*}
Prove that $C(p)$ is connected and closed in $X$.
\end{enumerate}

\vspace{1cm}
\noindent
Answer:\\
\\
(1)(a) $(X,\mathcal{U})$ a topological space. We are given that $C\subset X$ is connected.\\
Let $\overline{C} \supset C$ be the closure of $C$ and assume that it is not connected. Then $\exists\,U,V\in\mathcal{U}_{\overline{C}}$ both nonempty and disjoint such that:
\begin{align*}
U\cup V = \overline{C}.
\end{align*}
As $C$ is connected it is clear that $C$ itself is contained in $U$ or $V$ (say $U$).  If this were not the case
we would have found two open disjoint sets (i.e. $U \cap C$ and $V \cap C$) in $\mathcal{U}_{C}$ that seperated $C$.
Consequently, we must have that:
\begin{align*}
C\subset U \quad\text{and}\quad \overline{C}\cap V\neq\emptyset.
\end{align*}
As $\overline{C}\cap V\neq\emptyset$ let $p\in\overline{C}\cap V$, i.e. $p\in\overline{C}\setminus C$ (i.e. an accumulation point of $C$). But, by definition, if $p\in\overline{C}\setminus C$
\begin{align*}
\Rightarrow\quad \forall\,O\in\mathcal{U}\ \text{s.t.}\ p\in O
\quad\implies\quad O\cap C\neq\emptyset.
\end{align*}
But $V$ is one such $O \cap \overline{C}$ and hence:
\begin{align*}
\emptyset \neq V\cap C \subset V\cap U 
\end{align*}
\begin{center}
\underline{contradiction}.
\end{center}
Hence $C$ connected $\Rightarrow$ $\overline{C}$ also connected.\\
\qed\\
\\
(b) \noindent Assume $C_1,C_2\subset X$ are both connected.
Given $C_1\cap C_2\neq\emptyset$ show $C_1\cup C_2$ connected.

\noindent Assume $C_1\cup C_2$ is not connected. Then
$\exists\,U,V\in\mathcal{U}_{C_1\cup C_2}$ s.t. $U,V$ are disjoint and non-empty,
and
\begin{align*}
U\cup V = C_1\cup C_2.
\end{align*}
Clearly $C_1$ \& $C_2$ -- both connected -- must reside in different
$(U,V)$. Say $C_1\subset U$ \& $C_2\subset V$ - use same argument as in part (a) to show this.\\
\\
But
\begin{align*}
C_1\cap C_2 &\subset U\cap V\\
\text{and}\quad C_1\cap C_2\neq\emptyset
&\Rightarrow U\cap V\neq\emptyset.
\end{align*}
\begin{center}
Contradiction
\end{center}
Hence $C_1\cup C_2$ is connected.\\
\qed\\
\\
(c) \noindent Let
\begin{align*}
C(p)=\bigcup_{\substack{p\in C\\C\text{ connected}}}C.
\end{align*}
Assume $C(p)$ is not connected. Then $\exists\,U,V\in\mathcal{U}_{C(p)}$ both open,
disjoint, non-empty and $C(p) =U\cup V$.\\
A similar argument to part (b) follows. As each $C$ in $C(p)$
is connected it will belong to either $U$ or $V$.  $U \text{ and } V$ are both nonempty by assumption.  
This means that both $U$ and $V$ will contain different connected $C$'s.
Consequently, we must have that: $p \in U$ and $p \in V \implies p \in U \cap V$.
\begin{center}
contradiction.
\end{center}
Hence $C(p)$ is connected.\\
\\
\noindent Now $C(p)$ connected $\Rightarrow\overline{C(p)}$ connected (in part (1)).
\begin{align*}
&\Rightarrow C(p)\subset\overline{C(p)}\subset C(p)
\quad\text{from the definition of $C(p)$ (1.1).}\\
&\Rightarrow C(p)=\overline{C(p)}\\
&\Rightarrow C(p)\text{ is closed.}
\end{align*}
\qed

\newpage
\noindent\textbf{Exercise 2.}\\
\\
A topological space $(X,\mathcal{U})$ is called $T_1$ if for every pair of distinct points $x,y\in X$, there exists an open set $U$ such that
\begin{align*}
x\in U,\qquad y\notin U,
\end{align*}
and an open set $V$ such that
\begin{align*}
y\in V,\qquad x\notin V.
\end{align*}
Notice that the open sets $U$ and $V$ are not required to be disjoint.
\begin{enumerate}
\item[(a)] Prove that a topological space $(X,\mathcal{U})$ is $T_1$ if and only if every singleton $\{x\}$ is closed in $X$.
\end{enumerate}
Consider on $\mathbb{R}$ the topology of right rays
\begin{align*}
\mathcal{T}=\{(a,+\infty):a\in\mathbb{R}\}\cup\{\emptyset,\mathbb{R}\}.
\end{align*}
\begin{enumerate}
\item[(b)] Is $(\mathbb{R},\mathcal{T})$ Hausdorff? Is $(\mathbb{R},\mathcal{T})$ a $T_1$ space? Justify your answers.
\end{enumerate}


\vspace{1cm}
\noindent
Answer:\\
\\
(a) Let $(X,\mathcal{U})$ be a topological space.\\
(i) Assume every singleton set $\{x\}$ $(x\in X)$ is closed in $X$.

\noindent Given $x,y$ $(x\neq y)$ in $X$. Then:
\begin{align*}
U=X\setminus\{x\} &\quad\text{is open i.e. } U\in\mathcal{U}\\
V=X\setminus\{y\} &\quad\text{is open i.e. } V\in\mathcal{U}.
\end{align*}
Clearly:
\begin{align*}
&y\in U \quad\text{and}\quad x\notin U.\\
\text{\&}\quad &x\in V \quad\text{and}\quad y\notin V.
\end{align*}
Hence $(X,\mathcal{U})$ is $T_1$.\\
\\
(ii) Assume $(X,\mathcal{U})$ is $T_1$.\\
Given $x\in X$, then $\forall\,y\in X\setminus\{x\}$ - since $(X,\mathcal{U})$ is T1 - there is an open set $U_y$\\
such that: $y\in U_y$ \& $x\notin U_y$.\\
\\
Define:
\begin{align*}
U=\bigcup_{y\in X\setminus\{x\}}U_y.
\end{align*}
Then $U$ is also open (i.e. $U\in\mathcal{U}$).

\noindent Then $X\setminus U=\{x\}$ is closed.\\
\qed

\newpage
\noindent
(b) 
\begin{align*}
\mathcal{T}=\{(a,+\infty)\mid a\in\mathbb{R}\}\cup\{\emptyset,\mathbb{R}\}
\end{align*}
is a topology on $\mathbb{R}$.\\
\\
To determine if $\mathcal{T}$ is Hausdorff we must check that if given any $x,y\in\mathbb{R}$ $(x<y)$ there are
disjoint open sets that contain $x$ and $y$.\\
All open sets that contain $x$ are of the form $(p,+\infty)$ where $p<x$ - or $\mathbb{R}$ itself.\\
All open sets containing $y$ are of the form $(q,+\infty)$ where $q<y$ - or $\mathbb{R}$ itself.\\
As $x<y$. Any open set containing $x$ \underline{must} also contain $y$
$\Rightarrow (\mathbb{R},\mathcal{T})$ is \underline{not} Hausdorff.\\
\\
The same argument shows that \underline{any} open set that contains $x$ \underline{must} contain $y$.
So $(\mathbb{R},\mathcal{T})$ is also \underline{not} $T_1$.

\newpage
\noindent\textbf{Exercise 3.}

\noindent A map $f:X\to Y$ between topological spaces is called open if, for every open set $U\subseteq X$, $f(U)$ is open in $Y$.
\begin{enumerate}
\item[(a)] Given $(X,\mathcal{U})$ and $(Y,\mathcal{V})$ topological spaces, show that, if $f:X\to Y$ is bijective, continuous and open, then $f$ is a homeomorphism.
\end{enumerate}
Consider the torus $T\subseteq\mathbb{R}^3$, endowed with the topology $\mathcal{U}_T$ induced by the Euclidean topology on $\mathbb{R}^3$. Consider also $S^1\times S^1$ with the product topology.

\noindent Recall from class the following map, given in angular coordinates,
\begin{align*}
F:S^1\times S^1&\longrightarrow T,\\
(\theta,\varphi)&\longmapsto
\bigl((2+\cos(2\pi\theta))\cos(2\pi\varphi),
(2+\cos(2\pi\theta))\sin(2\pi\varphi),
\sin(2\pi\theta)\bigr).
\end{align*}
\begin{enumerate}
\item[(b)] Prove that $F$ is a homeomorphism.
\end{enumerate}
\newpage
\noindent\textbf{Exercise 4.} Let $(X,\mathcal{U})$ and $(Y,\mathcal{V})$ be non-empty topological spaces and consider $X\times Y$ endowed with the product topology. Prove that $X\times Y$ is connected if and only if both $X$ and $Y$ are connected.

\vspace{1cm}
\noindent
Answer:\\
\\
(i) Given that $(X \times Y, \mathcal{U}_{X \times Y})$ is connected prove that $(X,\mathcal{U}_{X})$ and $(Y,\mathcal{U}_{Y})$ are both connected.\\
\\
First we construct a map:
\begin{align*}
&f: (X \times Y, \mathcal{U}_{X \times Y}) \to (X,\mathcal{U}_{X}) \\
&f(x,y) = x
\end{align*}
This map is continuous since, for $O \in \mathcal{U}_{X} \implies f^{-1}(O) = O \times Y$ - which is open in $\mathcal{U}_{X \times Y}$. This map is also clearly surjective.  So, given that $(X \times Y, \mathcal{U}_{X \times Y})$ is connected we must have that $(X,\mathcal{U}_{X})$ is also connected.\\
\\
The same argument can be applied to show that $(Y,\mathcal{U}_{Y})$ must also be connected.\\ 
\\
(ii) Given that $(X,\mathcal{U})$ and $(Y,\mathcal{V})$ are both connected.\\
\\
Suppose $X \times Y = U \cup V$ is a separation into non-empty disjoint open sets. Choose $(x_0,y_0) \in U$.
The connected slice $\{x_0\} \times Y$ lies entirely in $U$. For every $y \in Y$, the connected slice $X \times \{y\}$ intersects it at $(x_0,y)$, so also lies entirely in $U$.\\
\\
Thus $U = X \times Y$, contradicting $V \neq \varnothing$.
\newpage
\noindent\textbf{Exercise 5.} Let $(X,\mathcal{U})$ and $(Y,\mathcal{V})$ be non-empty topological spaces and consider $X\times Y$ endowed with the product topology $\mathcal{U}_{X\times Y}$. Prove that $(X\times Y,\mathcal{U}_{X\times Y})$ is Hausdorff if and only if both $(X,\mathcal{U})$ and $(Y,\mathcal{V})$ are Hausdorff.

\vspace{1cm}
\noindent
Answer:\\
\\
(a) Assume $(X,\mathcal{U}_x),(Y,\mathcal{U}_y)$ are both Hausdorff.
Let $(x_1,y_1)$ \& $(x_2,y_2)$ be two distinct points in $X\times Y$.\\
(i) Assume $x_1\neq x_2$ \& $y_1\neq y_2$.
Then
\begin{align*}
&\exists\,U_1,U_2\in\mathcal{U}_x\ \text{s.t.}\ x_1\in U_1,\ x_2\in U_2,\quad U_1\cap U_2=\emptyset\\
\text{\&}\quad &\exists\,V_1,V_2\in\mathcal{U}_y\ \text{s.t.}\ y_1\in V_1,\ y_2\in V_2,\quad V_1\cap V_2=\emptyset.
\end{align*}
Then
\begin{align*}
&(x_1,y_1)\in U_1\times V_1;\quad (x_2,y_2)\in U_2\times V_2\\
\text{\&}\quad&(U_1\times V_1)\cap(U_2\times V_2)=\emptyset.
\end{align*}
Hence: $(X\times Y,\mathcal{U}_{X\times Y})$ is Hausdorff.\\
\\
(ii) If $y_1=y_2=y$ we just take $V\in\mathcal{U}_y$ s.t. $y\in V$.
Again
\begin{align*}
&(x_1,y)\in U_1\times V\quad\text{\&}\quad(x_2,y)\in U_2\times V\\
\text{\&}\quad &(U_1\times V)\cap(U_2\times V)=\emptyset.
\end{align*}
Hence: $(X\times Y,\mathcal{U}_{X\times Y})$ is again Hausdorff.\\
\\
(Same logic if $x_1=x_2$ \& $y_1\neq y_2$).
\begin{align*}
\therefore\quad(X\times Y,\mathcal{U}_{X\times Y})\text{ is Hausdorff}
\end{align*}
\\
(b) Assume that $(X\times Y,\mathcal{U}_{X\times Y})$ is Hausdorff.\\
Take any two distinct points $(x_1, y), (x_2, y) \in X\times Y$.  As $X\times Y$ Hausdorff there must exist
two open and disjoint "rectangles" $R_1 \; \& \; R_2$ s.t. $(x_i, y) \in R_i$.\\
\\
Let $R_1 = U_1 \times V_1$ and $R_2 = U_2 \times V_2$ - where $U_i \in \mathcal{U}_{X}$ and  $V_i \in \mathcal{U}_{Y}$.\\
Let $V = V_1 \cap V_2$, then $U_1 \times V$ and $U_2 \times V$ are also open disjoint "rectangles" in $X \times Y$ that seperate the
two points $(x_1, y)$ and $(x_2, y)$ - and their being disjoint means that $U_1 \cap U_2 = \emptyset$.\\
We conclude that $(X, \mathcal{U}_{X})$ is Hausdorf.\\
\\
The exact same argument can be applied for $(Y, \mathcal{U}_{Y})$ to show it is also Hausdorff.\\
\\
Hence we must have that both $(X, \mathcal{U}_{X})$ and $(Y, \mathcal{U}_{Y})$ are Hausdorff.
\qed
\newpage
\noindent\textbf{Exercise 6.} \textit{(Universal property of the Product Topology)}

\noindent Let $(X_\alpha,\mathcal{U}_\alpha)_{\alpha\in A}$ be a family of topological spaces. Consider the set $X:=\prod_{\alpha\in A}X_\alpha$ and, for every $\alpha\in A$, the projection $p_\alpha:X\to X_\alpha$. Prove that the product topology $\mathcal{U}_X$ on $X$ is the unique topology of $X$ such that, for every topological space $(Y,\mathcal{V})$ and map $f:Y\to X$, $f$ is continuous if and only if $p_\alpha\circ f:Y\to X_\alpha$ is continuous for every $\alpha\in A$.
\begin{align*}
\begin{array}{ccc}
(Y,\mathcal{V}) & \xrightarrow{\quad f\quad} & (X,\mathcal{U}_X)\\
& \underset{p_\alpha\circ f}{\searrow} & \big\downarrow{\scriptstyle p_\alpha}\\
&& (X_\alpha,\mathcal{U}_\alpha)
\end{array}
\end{align*}

\vspace{1cm}
\noindent
Answer:\\
\\


\end{document}
